\documentclass[8pt,a4paper,landscape]{extarticle}

\usepackage{fontspec}
\usepackage[russian]{babel}
\usepackage{amsmath,amssymb}
\usepackage{geometry}
\usepackage{microtype}
\usepackage{enumitem}
\usepackage{titlesec}
\usepackage{setspace}
\usepackage{parskip}
\usepackage{fancyvrb}
\usepackage{fvextra}
\usepackage{xcolor}
\usepackage{hyperref}
\usepackage{multicol}

\setmainfont{TeX Gyre Termes}
\setsansfont{TeX Gyre Heros}
\setmonofont{Latin Modern Mono}[Scale=MatchLowercase]

\geometry{
  left=5mm,
  right=5mm,
  top=5mm,
  bottom=6mm
}

\hypersetup{
  colorlinks=true,
  linkcolor=black,
  urlcolor=black
}

\pagestyle{empty}
\setstretch{0.92}
\setlength{\parindent}{0pt}
\setlength{\parskip}{1.2pt plus 0.3pt minus 0.2pt}
\setlength{\columnsep}{5mm}
\setlength{\columnseprule}{0.15pt}
\setlength{\abovedisplayskip}{2pt plus 0.5pt minus 0.5pt}
\setlength{\belowdisplayskip}{2pt plus 0.5pt minus 0.5pt}
\setlength{\abovedisplayshortskip}{1pt plus 0.5pt minus 0.5pt}
\setlength{\belowdisplayshortskip}{1pt plus 0.5pt minus 0.5pt}
\setlength{\jot}{0.5pt}

\titleformat{\section}
  {\normalfont\bfseries\normalsize}
  {\thesection}
  {0.35em}
  {}
\titlespacing*{\section}{0pt}{3pt plus 1pt minus 1pt}{1pt}

\titleformat{\subsection}
  {\normalfont\bfseries\small}
  {\thesubsection}
  {0.35em}
  {}
\titlespacing*{\subsection}{0pt}{3pt plus 1pt minus 1pt}{1pt}

\titleformat{\subsubsection}
  {\normalfont\bfseries\footnotesize}
  {\thesubsubsection}
  {0.35em}
  {}
\titlespacing*{\subsubsection}{0pt}{2pt plus 0.5pt minus 0.5pt}{1pt}

\setlist{
  nosep,
  leftmargin=*,
  topsep=0.5pt,
  partopsep=0pt,
  parsep=0pt,
  itemsep=0.5pt
}
\setlist[enumerate]{label=\arabic*.}

\DefineVerbatimEnvironment{Highlighting}{Verbatim}{
  breaklines,
  breaksymbolleft={},
  fontsize=\scriptsize,
  commandchars=\\\{\}
}
\DefineVerbatimEnvironment{verbatim}{Verbatim}{
  breaklines,
  breaksymbolleft={},
  fontsize=\scriptsize
}

\newcommand{\tightlist}{}
\sloppy

\begin{document}
\begin{multicols*}{4}

\subsection{Блок 1}\label{ux431ux43bux43eux43a-1}

\subsubsection{Задача 5 (Определение 2-раундовой сети Фейстеля)}\label{ux437ux430ux434ux430ux447ux430-5-ux43eux43fux440ux435ux434ux435ux43bux435ux43dux438ux435-2-ux440ux430ux443ux43dux434ux43eux432ux43eux439-ux441ux435ux442ux438-ux444ux435ux439ux441ux442ux435ux43bux44f}

\textbf{Решение:} Двухраундовая сеть Фейстеля не является псевдослучайной перестановкой (PRP), так как она не обеспечивает полную диффузию. Это позволяет построить различитель (distinguisher).

\begin{enumerate}
\def\labelenumi{\arabic{enumi}.}
\item
  \textbf{Анализ схемы:} Пусть входной 64-битный блок делится на две 32-битные половины: левую \(L_0\) и правую \(R_0\). (На схеме \(R_0\) сверху, \(L_0\) снизу, но логика перекрестий стандартная). Опишем преобразования по раундам:

  \begin{itemize}
  \tightlist
  \item
    \textbf{Раунд 1:} \(L_1 = R_0\) \(R_1 = L_0 \oplus F(k_1, R_0)\)
  \item
    \textbf{Раунд 2:} \(L_2 = R_1 = L_0 \oplus F(k_1, R_0)\) \(R_2 = L_1 \oplus F(k_2, R_1) = R_0 \oplus F(k_2, R_1)\) Выход сети состоит из половин \(L_2\) и \(R_2\) (обычно конкатенируются как \(L_2 \parallel R_2\)).
  \end{itemize}
\item
  \textbf{Применение заданных входов:} Нам даны два входа:

  \begin{itemize}
  \tightlist
  \item
    Вход A: \(0^{64}\). То есть \(L_{0A} = 0^{32}\), \(R_{0A} = 0^{32}\).
  \item
    Вход B: \(1^{32}0^{32}\). То есть \(L_{0B} = 1^{32}\) (все единицы), \(R_{0B} = 0^{32}\).
  \end{itemize}

  Заметим ключевую особенность: \textbf{правые половины обоих входов равны} (\(R_{0A} = R_{0B} = 0^{32}\)).
\item
  \textbf{Поиск коллизии/закономерности:} Посмотрим, что произойдет с половиной \(L_2\) для обоих входов:

  \begin{itemize}
  \tightlist
  \item
    Для входа A: \(L_{2A} = L_{0A} \oplus F(k_1, R_{0A}) = 0^{32} \oplus F(k_1, 0^{32}) = F(k_1, 0^{32})\)
  \item
    Для входа B: \(L_{2B} = L_{0B} \oplus F(k_1, R_{0B}) = 1^{32} \oplus F(k_1, 0^{32})\)
  \end{itemize}

  Сложим их по модулю 2 (операция XOR, \(\oplus\)): \(L_{2A} \oplus L_{2B} = F(k_1, 0^{32}) \oplus 1^{32} \oplus F(k_1, 0^{32}) = \mathbf{1^{32}}\)

  \textbf{Вывод:} Если выходы сгенерированы этой 2-раундовой сетью Фейстеля, то одна из 32-битных половин (левая или правая, в зависимости от порядка конкатенации) у двух шифротекстов при операции XOR должна дать \(1^{32}\). В шестнадцатеричном виде \(1^{32}\) --- это строка из восьми букв \texttt{f} (\texttt{ffffffff}). Истинная случайная перестановка даст такой результат с пренебрежимо малой вероятностью (\(1/2^{32}\)).
\item
  \textbf{Проверка вариантов ответа:} Ищем вариант, где одна из половин первого выхода при XOR со второй половиной дает \texttt{ffffffff}. (В 16-ричной системе \(x \oplus y = f \implies\) биты полностью инвертированы).

  \begin{itemize}
  \item
    \textbf{Вариант 1:} Половина 1: \texttt{7b50baab} \(\oplus\) \texttt{ac343a22} \(\neq\) \texttt{ffffffff} (т.к. 7 \(\oplus\) a = d)
  \item
    \textbf{Вариант 2:} Половина 1: \texttt{5f67abaf} \(\oplus\) \texttt{bbe033c0} \(\neq\) \texttt{ffffffff} (т.к. 5 \(\oplus\) b = e)
  \item
    \textbf{Вариант 3:} Половина 1: \texttt{290b6e3a} \(\oplus\) \texttt{d6f491c5}. Проверим посимвольно: 2 \(\oplus\) d = 0010 \(\oplus\) 1101 = 1111 (f) 9 \(\oplus\) 6 = 1001 \(\oplus\) 0110 = 1111 (f) 0 \(\oplus\) f = 0000 \(\oplus\) 1111 = 1111 (f) b \(\oplus\) 4 = 1011 \(\oplus\) 0100 = 1111 (f) 6 \(\oplus\) 9 = 0110 \(\oplus\) 1001 = 1111 (f) e \(\oplus\) 1 = 1110 \(\oplus\) 0001 = 1111 (f) 3 \(\oplus\) c = 0011 \(\oplus\) 1100 = 1111 (f) a \(\oplus\) 5 = 1010 \(\oplus\) 0101 = 1111 (f) \textbf{Результат: \texttt{ffffffff}. Идеальное совпадение.}
  \item
    \textbf{Вариант 4} (для убедительности): Половина 1: \texttt{9d1a4f78} \(\oplus\) \texttt{75e5e3ea} \(\neq\) \texttt{ffffffff} (т.к. 9 \(\oplus\) 7 = e)
  \end{itemize}
\end{enumerate}

\textbf{Ответ: 3} (Строки из 3-го варианта являются выходом данной сети Фейстеля, так как XOR их первых 32-битных блоков равен \texttt{ffffffff}, что доказывает структурную уязвимость 2-х раундов).

\subsection{Блок 2}\label{ux431ux43bux43eux43a-2}

\subsubsection{Задача 1}\label{ux437ux430ux434ux430ux447ux430-1}

\textbf{Решение:} В режиме счетчика (CTR) шифрование и расшифрование идет наложением гаммы (XOR): Шифрование: \(C_i = P_i \oplus E_K(CTR_i)\) Расшифрование: \(P_i = C_i \oplus E_K(CTR_i)\)

Генерация гаммы \(E_K(CTR_i)\) зависит \textbf{только} от ключа и значения счетчика. Она никак не зависит от предыдущих или текущих блоков шифротекста \(C\). Т.к. ошибка (искажение битов) произошла только в блоке \(C_{L/2}\), то при операции \(P_{L/2} = C_{L/2} \oplus E_K(CTR_{L/2})\) испортится только один этот блок открытого текста. Остальная гамма сгенерируется корректно.

\textbf{Ответ: 1}

\begin{center}\rule{0.5\linewidth}{0.5pt}\end{center}

\subsubsection{Задача 2}\label{ux437ux430ux434ux430ux447ux430-2}

\textbf{Решение:} Из условия мы знаем, как алгоритм вычисляет тэг. Злоумышленник сделал запрос для сообщения \(m\) со случайным \(IV = r\) и получил тэг \(t\). Формула тэга: \(t = E_K(m \oplus r)\).

Чтобы совершить подделку (forgery), нам нужно сгенерировать такой новый вектор \(IV'\) и новое сообщение \(m'\), чтобы при вычислении \(E_K(m' \oplus IV')\) под шифрование попало то же самое значение, что и в первый раз (тогда мы сможем использовать старый тэг \(t\)). То есть нам нужно, чтобы выполнялось равенство: \(m' \oplus IV' = m \oplus r\)

Проверим \textbf{четвертый} вариант ответа из списка: Предлагается тэг \textbf{\((r \oplus 1^n, t)\)} для нового сообщения \textbf{\(m \oplus 1^n\)}. Здесь: Новый вектор \(IV' = r \oplus 1^n\) Новое сообщение \(m' = m \oplus 1^n\)

Подставим их в проверку алгоритма верификации: \(m' \oplus IV' = (m \oplus 1^n) \oplus (r \oplus 1^n)\)

Так как операция XOR ассоциативна и коммутативна, перегруппируем: \(= m \oplus r \oplus (1^n \oplus 1^n)\)

Одинаковые значения при сложении по модулю 2 дают ноль (\(1^n \oplus 1^n = 0^n\)): \(= m \oplus r \oplus 0^n = m \oplus r\)

Получается, что для этого нового сообщения и нового вектора \(IV\) алгоритм будет вычислять \(E_K(m \oplus r)\). Но мы уже знаем из первого перехвата, что \(E_K(m \oplus r) = t\). Значит, старый тэг \(t\) идеально подойдет! Подделка успешна.

\emph{(Для справки, вариант 3 с \texttt{(m\ ⊕\ t,\ t)} дал бы \(E_K(0^n \oplus m \oplus t) = E_K(m \oplus t) \neq t\), поэтому он неверен).}

\textbf{Ответ:} Выбираем вариант \textbf{- Тег \((r \oplus 1^n, t)\) -- валидный для 1-блочного сообщения \(m \oplus 1^n\)}

\begin{center}\rule{0.5\linewidth}{0.5pt}\end{center}

\subsubsection{Задача 3}\label{ux437ux430ux434ux430ux447ux430-3}

\textbf{Решение:} RSA: \(N=33, e=3, c=3\). Найти \(m = c^d \pmod N\).

\begin{enumerate}
\def\labelenumi{\arabic{enumi}.}
\tightlist
\item
  Факторизация \(N\): \(N = 33 = 3 \times 11 \implies p=3, q=11\).
\item
  Функция Эйлера: \(\phi(N) = (p-1)(q-1) = 2 \times 10 = 20\).
\item
  Поиск секретной экспоненты \(d\): \(e \cdot d \equiv 1 \pmod{\phi(N)}\) \(3 \cdot d \equiv 1 \pmod{20}\) \(3 \cdot 7 = 21 \equiv 1 \pmod{20} \implies d = 7\).
\item
  Расшифрование: \(m = 3^7 \pmod{33}\) Считаем: \(3^3 = 27 \equiv -6 \pmod{33}\) \(3^6 = (-6)^2 = 36 \equiv 3 \pmod{33}\) \(3^7 = 3^6 \cdot 3 \equiv 3 \cdot 3 = 9 \pmod{33}\)
\end{enumerate}

\textbf{Ответ: m = 9}

\begin{center}\rule{0.5\linewidth}{0.5pt}\end{center}

\subsubsection{Задача 4}\label{ux437ux430ux434ux430ux447ux430-4}

\textbf{Решение:} Проверяем свойства \(H'\) при условии, что \(H\) криптографически стойкая (CRHF):

\begin{enumerate}
\def\labelenumi{\arabic{enumi}.}
\tightlist
\item
  \(H'(m) = H(m) \oplus H(m \oplus 1^{|m|})\) --- \textbf{Нет}. Коллизия тривиальна: \(H'(x) = H'(\bar{x})\) в силу коммутативности XOR.
\item
  \(H'(m) = H(m)[0\dots31]\) --- \textbf{Нет}. Выход всего 32 бита. По парадоксу дней рождений коллизия ищется за \(\approx 2^{16}\) операций (небезопасно).
\item
  \(H'(m) = H(H(m))\) --- \textbf{Да, устойчива}. \emph{Доказательство от противного:} пусть нашли \(m_1 \neq m_2\), что \(H(H(m_1)) = H(H(m_2))\). Это значит, либо \(H(m_1) = H(m_2)\) (нашли коллизию для \(H\), противоречие), либо \(X = H(m_1), Y = H(m_2)\) и \(H(X) = H(Y)\) при \(X \neq Y\) (снова нашли коллизию для \(H\), противоречие).
\item
  \(H'(m) = H(m) \oplus H(m)\) --- \textbf{Нет}. Всегда выдает 0. \(H'(x) = H'(y) = 0\).
\item
  \(H'(m) = H(0)\) --- \textbf{Нет}. Константа, коллизии для любых \(x, y\).
\item
  \(H'(m) = H(m) \parallel H(m)\) --- \textbf{Да, устойчива}. Если конкатенации равны, то равны и части \(\implies H(x)=H(y) \implies\) сводится к стойкости \(H\).
\end{enumerate}

\textbf{Ответ:} Основная устойчивая функция: \textbf{\(H'(m) = H(H(m))\)}. \emph{(Также формально подходит \(H'(m) = H(m) \parallel H(m)\))}.

\begin{center}\rule{0.5\linewidth}{0.5pt}\end{center}

\subsubsection{Задача 5}\label{ux437ux430ux434ux430ux447ux430-5}

\textbf{Решение:} Схема CBC: \(IV = E_K(nonce), C_i = E_K(m_i \oplus C_{i-1})\). Обозначим \(E_K\) как \(E\).

\textbf{Запрос 1:} \(m = (0^\ell, 0^\ell)\), \(nonce = 0^\ell\). \(IV_1 = E(0^\ell)\) \(c_0 = E(m_0 \oplus IV_1) = E(0^\ell \oplus E(0^\ell)) = E(E(0^\ell))\) \(c_1 = E(m_1 \oplus c_0) = E(0^\ell \oplus c_0) = E(c_0)\) \emph{Вывод (подсказка):} \textbf{\(c_1 = E(c_0)\)}

\textbf{Запрос 2:} \(m_1 = c_0 \oplus c_1\), \(nonce = c_0\). Новый \(IV_2 = E(nonce) = E(c_0)\). Т.к. из первого запроса мы знаем, что \(E(c_0) = c_1\), след-но: \textbf{\(IV_2 = c_1\)}. Вычисляем шифротекст: \(c_0' = E(m_1 \oplus IV_2)\) Подставляем \(m_1\) и \(IV_2\): \(c_0' = E((c_0 \oplus c_1) \oplus c_1) = E(c_0 \oplus 0^\ell) = E(c_0)\) Снова используем вывод из первого запроса (\(E(c_0) = c_1\)): \textbf{\(c_0' = c_1\)}

\textbf{Ответ:} Выбираем вариант \textbf{\(c_1 = c_0'\)}

\subsection{Блок 3}\label{ux431ux43bux43eux43a-3}

\subsubsection{Задача 1. Ошибка в режиме CBC}\label{ux437ux430ux434ux430ux447ux430-1.-ux43eux448ux438ux431ux43aux430-ux432-ux440ux435ux436ux438ux43cux435-cbc}

\textbf{Решение:} Формула расшифрования в режиме CBC (Cipher Block Chaining): \(P_i = AES^{-1}_K(C_i) \oplus C_{i-1}\)

Если при передаче искажен (содержит ошибки) блок шифротекста \(C_{L/2}\), проанализируем, на какие блоки открытого текста это повлияет:

\begin{enumerate}
\def\labelenumi{\arabic{enumi}.}
\tightlist
\item
  \textbf{\(P_{L/2}\)}: Вычисляется как \(AES^{-1}_K(C_{L/2}) \oplus C_{L/2 - 1}\). Так как \(C_{L/2}\) поступает на вход алгоритму дешифрования, из-за лавинного эффекта блочного шифра блок \(P_{L/2}\) будет искажен полностью (псевдослучайный мусор).
\item
  \textbf{\(P_{L/2 + 1}\)}: Вычисляется как \(AES^{-1}_K(C_{L/2 + 1}) \oplus C_{L/2}\). Блок \(C_{L/2}\) здесь используется только в операции XOR. Следовательно, в блоке \(P_{L/2 + 1}\) инвертируются ровно те биты, которые были повреждены в \(C_{L/2}\). Блок испорчен.
\item
  \textbf{\(P_{L/2 + 2}\)}: Вычисляется как \(AES^{-1}_K(C_{L/2 + 2}) \oplus C_{L/2 + 1}\). Оба нужных блока шифротекста целые. Расшифровка корректна.
\end{enumerate}

Дальнейшее распространение ошибки не происходит. Испорчено ровно 2 блока.

\textbf{Ответ: 2}

\begin{center}\rule{0.5\linewidth}{0.5pt}\end{center}

\subsubsection{Задача 2. Длина сообщения и Padding в CBC}\label{ux437ux430ux434ux430ux447ux430-2.-ux434ux43bux438ux43dux430-ux441ux43eux43eux431ux449ux435ux43dux438ux44f-ux438-padding-ux432-cbc}

\textbf{Решение:} Проанализируем размер пакета. Общий размер пакета = 128 байт. Размер блока AES и вектора инициализации (IV) = 16 байт. Полезная нагрузка (Ciphertext, CT) = \(128 - 16 (IV) = 112\) байт. \(112\) байт составляют ровно \(112 / 16 = 7\) блоков AES.

Режим CBC для произвольной длины требует применения схемы дополнения (padding), обычно PKCS\#7. Дополнение всегда добавляет от 1 до 16 байт (даже если исходное сообщение кратно размеру блока, добавляется целый блок паддинга). Следовательно, размер исходного открытого текста \(L_{msg}\) должен удовлетворять условию: \(CT\_length - 16 \le L_{msg} < CT\_length\) \(112 - 16 \le L_{msg} \le 112 - 1 \implies\) \textbf{Длина текста должна быть от 96 до 111 символов.}

Оценим длину предложенных вариантов (включая пробелы и знаки препинания):

\begin{enumerate}
\def\labelenumi{\arabic{enumi}.}
\tightlist
\item
  \emph{The significance\ldots{}} --- явно больше 150 символов (не влезает).
\item
  \emph{If qualified opinions\ldots{}} --- 125 символов (не влезает в 112 байт шифротекста).
\item
  \textbf{In this letter I make some remarks on a general principle relevant to enciphering in general and my machine.} --- ровно \textbf{108 символов}. Попадает в диапазон {[}96, 111{]}. Потребует 4 байта паддинга (0x04 04 04 04).
\item
  \emph{The most direct computation\ldots{}} --- 92 символа. (Если текст 92 байта, паддинг достроит его до \(92+4 = 96\) байт = 6 блоков. Размер пакета с IV был бы \(96+16 = 112\) байт, а у нас по условию 128 байт).
\end{enumerate}

\textbf{Ответ: 3-е сообщение} (`In this letter I make some remarks on a general principle relevant to enciphering in general and my machine.')

\begin{center}\rule{0.5\linewidth}{0.5pt}\end{center}

\subsubsection{Задача 3. RSA}\label{ux437ux430ux434ux430ux447ux430-3.-rsa}

\textbf{Решение:} Дано: \(N=35, e=7, c=3\). Найти \(m\).

\begin{enumerate}
\def\labelenumi{\arabic{enumi}.}
\tightlist
\item
  Факторизуем модуль: \(N = 35 = 5 \times 7 \implies p=5, q=7\).
\item
  Вычисляем функцию Эйлера: \(\phi(N) = (p-1)(q-1) = 4 \times 6 = 24\).
\item
  Находим приватную экспоненту \(d\), такую что \(e \cdot d \equiv 1 \pmod{\phi(N)}\): \(7 \cdot d \equiv 1 \pmod{24}\). Заметим, что \(7 \cdot 7 = 49 = 24 \cdot 2 + 1 \implies d = 7\).
\item
  Дешифруем сообщение: \(m = c^d \pmod N = 3^7 \pmod{35}\). Считаем: \(3^2 = 9\) \(3^3 = 27\) \(3^4 = 81 \equiv 11 \pmod{35}\) (т.к. \(35 \times 2 = 70\)) \(3^6 = 3^4 \cdot 3^2 \equiv 11 \cdot 9 = 99 \equiv -6 \pmod{35}\) (т.к. \(35 \times 3 = 105\)) \(3^7 = 3^6 \cdot 3 \equiv -6 \cdot 3 = -18 \pmod{35}\) Положительный остаток: \(35 - 18 = 17\).
\end{enumerate}

\textbf{Ответ: m = 17}

\begin{center}\rule{0.5\linewidth}{0.5pt}\end{center}

\subsubsection{Задача 4. Коллизия для функции сжатия (Davies-Meyer)}\label{ux437ux430ux434ux430ux447ux430-4.-ux43aux43eux43bux43bux438ux437ux438ux44f-ux434ux43bux44f-ux444ux443ux43dux43aux446ux438ux438-ux441ux436ux430ux442ux438ux44f-davies-meyer}

\textbf{Решение:} Функция сжатия задана как \(f_1(x,y) = AES(y,x) \oplus y\). \emph{Важно учесть нотацию из условия:} \(AES(x,y)\) означает ключ \(x\), сообщение \(y\). Следовательно, в записи \(AES(y,x)\) \textbf{ключом является \(y\)}, а \textbf{сообщением \(x\)}. Обозначим это как \(E_y(x)\). Тогда \(f_1(x,y) = E_y(x) \oplus y\).

Нам нужно найти две пары \((x_1, y_1)\) и \((x_2, y_2)\) такие, что: \(E_{y_1}(x_1) \oplus y_1 = E_{y_2}(x_2) \oplus y_2\) Пусть \(v = E_{y_1}(x_1)\) (как предложено во всех вариантах). Тогда левая часть равна \(v \oplus y_1\). Подставляем в уравнение: \(v \oplus y_1 = E_{y_2}(x_2) \oplus y_2\) \(E_{y_2}(x_2) = v \oplus y_1 \oplus y_2\) Так как мы сами выбираем ключ \(y_2\), мы можем просто применить операцию расшифрования \(AES^{-1}\) (где ключом выступает первый аргумент \(y_2\)) к правой части, чтобы найти \(x_2\): \(x_2 = AES^{-1}(y_2, v \oplus y_1 \oplus y_2)\)

Смотрим предложенные варианты. Ровно этот вывод содержится в третьем варианте ответа. (Варианты с \(y_2 = AES^{-1}(...)\) не имеют криптографического смысла, т.к. дешифратор ищет открытый текст, а не ключ).

\textbf{Ответ: Выбираем 3-й вариант:}

\begin{itemize}
\tightlist
\item
  Выбираем \(x_1, y_1, y_2\) произвольно (\(y_1 \neq y_2\)) и полагаем \(v := AES(y_1, x_1)\), \(x_2 := AES^{-1}(y_2, v \oplus y_1 \oplus y_2)\)
\end{itemize}

\begin{center}\rule{0.5\linewidth}{0.5pt}\end{center}

\subsubsection{Задача 5. Безопасные MAC}\label{ux437ux430ux434ux430ux447ux430-5.-ux431ux435ux437ux43eux43fux430ux441ux43dux44bux435-mac}

\textbf{Решение:} Обозначим исходный безопасный MAC как \(S\), а злоумышленника как \(\mathcal{A}\). Схема считается небезопасной, если \(\mathcal{A}\) может выполнить экзистенциальную подделку (existential forgery) --- найти валидный тэг для хотя бы одного нового сообщения, которое он ранее не запрашивал у оракула.

По условию исходный MAC определен на пространствах \(K, M, T\), где \(M=\{0,1\}^n\) и \(T=\{0,1\}^{128}\).

Анализируем предложенные конструкции:

\begin{enumerate}
\def\labelenumi{\arabic{enumi}.}
\item
  \textbf{\(S'(k, m) = [t \leftarrow S(k,m), \text{output } (t, t)]\)} \textbf{Безопасен.} Сводится к безопасности \(S\). Чтобы подделать тэг \((t^*, t^*)\) для нового \(m^*\), \(\mathcal{A}\) должен найти \(t^* = S(k, m^*)\). Без знания ключа это эквивалентно взлому базовой схемы \(S\), что по условию невозможно.
\item
  \textbf{\(S'(k, m) = S(k, m \oplus m)\)} \textbf{Небезопасен.} Так как \(m \oplus m = 0^n\) для любого \(m\), функция всегда выдает константный тэг \(t = S(k, 0^n)\). \emph{Атака:} \(\mathcal{A}\) запрашивает тэг для произвольного \(m_1\), получает \(t\). Затем заявляет \(t\) как валидный тэг для любого другого \(m_2 \neq m_1\). Проверка пройдет успешно.
\item
  \textbf{\(S'(k, m) = (S(k, m), S(k, 0^n))\)} \textbf{Небезопасен.} Возможна подделка для конкретного сообщения \(m^* = 0^n\). \emph{Атака:} \(\mathcal{A}\) запрашивает тэг для любого сообщения \(m \neq 0^n\). Оракул возвращает пару \((t_1, t_2)\), где \(t_2 = S(k, 0^n)\). Теперь \(\mathcal{A}\) нужно выдать тэг для нового сообщения \(m^* = 0^n\). Согласно алгоритму, правильный тэг для него --- это \((S(k, 0^n), S(k, 0^n))\). Но \(\mathcal{A}\) уже знает значение \(S(k, 0^n)\), это \(t_2\). \(\mathcal{A}\) выдает пару \((t_2, t_2)\) как тэг для \(0^n\). Подделка успешна.
\item
  \textbf{\(S'(k, m) = S(k, m \oplus 1^n)\)} \textbf{Безопасен.} Операция XOR с константой \(f(m) = m \oplus 1^n\) является биекцией в пространстве \(\{0,1\}^n\). Любой запрос на подделку нового сообщения \(m^*\) эквивалентен попытке подделать базовый MAC \(S\) для сообщения \((m^* \oplus 1^n)\), которое тоже является новым (т.к. функция биективна, \(\mathcal{A}\) ранее его не запрашивал). Сводится к безопасности \(S\).
\item
  \textbf{\(S'(k, m) = S(k, m || m)\)} \textbf{Формально некорректен / Безопасен (с оговоркой).} По условию домен базового MAC строго ограничен: \(M=\{0,1\}^n\). Сообщение \(m||m\) имеет длину \(2n\), поэтому функция \(S\) от него не определена. \emph{Если предположить}, что \(S\) принимает сообщения произвольной длины (или длины \(\ge 2n\)), то конструкция \textbf{безопасна}. Функция \(m \mapsto m||m\) инъективна. Чтобы подделать тэг для нового \(m^*\), нужно вычислить базовый MAC для строки \(m^* || m^*\), которая никогда ранее не поступала на вход оракулу.
\item
  \textbf{\(S'(k, m) = S(k, m[0\dots n-2] || 0)\)} \textbf{Небезопасен.} Принудительное зануление последнего бита порождает коллизии. \emph{Атака:} \(\mathcal{A}\) запрашивает тэг для сообщения, заканчивающегося нулем (например, \(0^n\)), получает \(t\). Затем выдает \(t\) в качестве тэга для сообщения, отличающегося только последним битом (например, \(0^{n-1}||1\)). При верификации последний бит занулится, и тэг совпадет. Подделка успешна.
\end{enumerate}

\textbf{Ответ:} Безопасны только варианты \textbf{1} и \textbf{4} (и вариант \textbf{5}, если расширить область определения базового MAC). Варианты 2, 3 и 6 небезопасны, так как допускают экзистенциальную подделку.

\subsection{Блок 4}\label{ux431ux43bux43eux43a-4}

\subsubsection{Задача 1. Оценка времени полного перебора AES-128}\label{ux437ux430ux434ux430ux447ux430-1.-ux43eux446ux435ux43dux43aux430-ux432ux440ux435ux43cux435ux43dux438-ux43fux43eux43bux43dux43eux433ux43e-ux43fux435ux440ux435ux431ux43eux440ux430-aes-128}

\textbf{Решение:}

\begin{enumerate}
\def\labelenumi{\arabic{enumi}.}
\tightlist
\item
  Размер пространства ключей AES-128 составляет \(2^{128}\) ключей.
\item
  Бюджет организации: \(\$4\ 000\ 000\ 000\ 000 = 4 \cdot 10^{12}\) долларов.
\item
  Стоимость одной машины: \(\$200\).
\item
  Количество машин, которые можно купить: \(N = (4 \cdot 10^{12}) / 200 = 2 \cdot 10^{10}\) машин.
\item
  Скорость одной машины: \(10^9\) ключей/сек.
\item
  Суммарная скорость перебора: \(V = 2 \cdot 10^{10} \cdot 10^9 = 2 \cdot 10^{19}\) ключей/сек.
\item
  Для гарантированного подбора (исчерпывающего поиска) нужно перебрать все \(2^{128}\) вариантов. Оценим \(2^{128}\) в десятичной системе (зная, что \(2^{10} \approx 10^3\)): \(2^{128} = 2^8 \cdot (2^{10})^{12} \approx 256 \cdot (10^3)^{12} = 256 \cdot 10^{36} = 2.56 \cdot 10^{38}\) ключей.
\item
  Время перебора в секундах: \(t = \frac{2.56 \cdot 10^{38}}{2 \cdot 10^{19}} = 1.28 \cdot 10^{19}\) секунд.
\item
  Переведем в года (в году \(\approx 3.15 \cdot 10^7\) секунд): \(t_{years} = \frac{1.28 \cdot 10^{19}}{3.15 \cdot 10^7} \approx 4.06 \cdot 10^{11}\) лет.
\end{enumerate}

\textbf{Ответ:} Потребуется приблизительно \textbf{\(4 \cdot 10^{11}\) лет} (около 400 миллиардов лет). Вычислительных мощностей за такие деньги всё ещё катастрофически не хватает для взлома AES-128.

\begin{center}\rule{0.5\linewidth}{0.5pt}\end{center}

\subsubsection{Задача 2. Построение безопасных PRF}\label{ux437ux430ux434ux430ux447ux430-2.-ux43fux43eux441ux442ux440ux43eux435ux43dux438ux435-ux431ux435ux437ux43eux43fux430ux441ux43dux44bux445-prf}

\textbf{Решение:} Исходная функция \(F(k, x)\) --- безопасная PRF. Проверяем конструкции:

\begin{enumerate}
\def\labelenumi{\arabic{enumi}.}
\tightlist
\item
  \(F'(k, x) = F(k, x) || 0\). \textbf{Небезопасна.} Выход всегда заканчивается нулем. Злоумышленник отличит её от истинной случайной функции \(f: \{0,1\}^n \to \{0,1\}^{n+1}\) по одному запросу (у случайной последний бит равен 1 с вероятностью 1/2).
\item
  \(F'((k_1, k_2), x) = F(k_1, x) || F(k_2, x)\). \textbf{Безопасна.} Так как ключи независимы, это конкатенация выходов двух независимых PRF. Результат неотличим от случайной функции с областью значений \(\{0,1\}^{2n}\).
\item
  \(F'(k, x) = reverse(F(k, x))\). \textbf{Безопасна.} Реверс (перестановка) битов --- это биективное преобразование. Если выход \(F\) неотличим от случайной строки, то и перевернутая строка неотличима от случайной.
\item
  \(F'(k, x) = F(k, x)\) если \(x \neq 0^n\), иначе \(0^n\). \textbf{Небезопасна.} Злоумышленник делает запрос \(x = 0^n\) и получает \(0^n\). У истинной случайной функции вероятность выдать ровно \(0^n\) равна \(1/2^n\) (ничтожно мала). Атака-различитель тривиальна.
\item
  \(F'(k, x) = F(k, x)\) если \(x \neq 0^n\), иначе \(k\). \textbf{Небезопасна.} При запросе \(x = 0^n\) злоумышленник получает ключ. Далее он может локально вычислять \(F(k, x)\) для любых \(x\) и со 100\% вероятностью отличать функцию от случайной.
\item
  \(F'(k, x) = F(k, x)[0 \dots n-2]\). \textbf{Безопасна.} Усечение выхода безопасной PRF сохраняет свойства псевдослучайности (выход неотличим от случайной функции в \(\{0,1\}^{n-1}\)).
\end{enumerate}

\textbf{Ответ: Безопасными являются конструкции 2, 3 и 6.}

\begin{center}\rule{0.5\linewidth}{0.5pt}\end{center}

\subsubsection{Задача 3. Уязвимости MAC в файловой системе}\label{ux437ux430ux434ux430ux447ux430-3.-ux443ux44fux437ux432ux438ux43cux43eux441ux442ux438-mac-ux432-ux444ux430ux439ux43bux43eux432ux43eux439-ux441ux438ux441ux442ux435ux43cux435}

\textbf{Решение:} По условию MAC применяется \textbf{исключительно к содержимому файла} (\(Tag = S(k, Content)\)). Любые манипуляции с самим содержимым (стирание байта, добавление данных) приведут к изменению \(Content\), и верификация MAC закономерно завершится ошибкой (атака предотвращена). Замена тега и файла с другого компьютера с \emph{другим ключом} также вызовет ошибку верификации, так как ключ локального компьютера не совпадет (атака предотвращена). А вот метаданные (включая дату модификации, имя файла, права доступа) не входят в вычисление MAC. Следовательно, их можно менять безнаказанно --- верификатор проверит только содержимое и скажет, что оно не искажено.

\textbf{Ответ:} Система не предотвращает атаку: \textbf{Изменение времени последней модификации файла.}

\begin{center}\rule{0.5\linewidth}{0.5pt}\end{center}

\subsubsection{Задача 4. Распределение ключей (по изображению)}\label{ux437ux430ux434ux430ux447ux430-4.-ux440ux430ux441ux43fux440ux435ux434ux435ux43bux435ux43dux438ux435-ux43aux43bux44eux447ux435ux439-ux43fux43e-ux438ux437ux43eux431ux440ux430ux436ux435ux43dux438ux44e}

\textbf{Решение:} По условию имеется 6 пользователей. Мы должны раздать им ключи так, чтобы ни одно множество ключей не было подмножеством другого. На изображении дан частичный список множеств: \(S_1 = \{k_2, k_4\}\), \(S_2 = \{k_2, k_3\}\), \(S_3 = \{k_3, k_4\}\), \(S_4 = \{k_1, k_3\}\)\ldots{} Видно, что каждому выдается ровно по 2 ключа. Чтобы обеспечить 6 уникальных наборов из пула ключей, не являющихся подмножествами друг друга, используется сочетание из \(N\) по 2. Нам нужно \(\binom{N}{2} \ge 6\). Минимальное \(N=4\), так как \(\binom{4}{2} = 6\). Следовательно, используется ровно 4 ключа (\(k_1, k_2, k_3, k_4\)), и 6 пользователей получают все возможные пары. Найдем недостающие множества \(S_5\) и \(S_6\), перебрав все пары: Уже есть: \(\{k_2, k_4\}, \{k_2, k_3\}, \{k_3, k_4\}, \{k_1, k_3\}\). Остались пары с участием \(k_1\). \textbf{Ответ:} Недостающие наборы ключей для пользователей B5 и B6: \textbf{\(S_5 = \{k_1, k_2\}\)} \textbf{\(S_6 = \{k_1, k_4\}\)}

\begin{center}\rule{0.5\linewidth}{0.5pt}\end{center}

\subsubsection{Задача 5. Пересчет ECBC-MAC при изменении бита}\label{ux437ux430ux434ux430ux447ux430-5.-ux43fux435ux440ux435ux441ux447ux435ux442-ecbc-mac-ux43fux440ux438-ux438ux437ux43cux435ux43dux435ux43dux438ux438-ux431ux438ux442ux430}

\textbf{Решение:} В схеме Encrypted CBC-MAC (ECBC-MAC) используется два ключа: \(k_1\) для основной цепи CBC и \(k_2\) для финального шифрования. Пусть \(c_{i}\) --- промежуточные блоки шифротекста в цепи. Тэг для исходного сообщения \(m\): \(c_n = E(k_1, m_n \oplus c_{n-1})\) \(Tag(m) = E(k_2, c_n)\)

Нам дан \(Tag(m)\), ключи \(k_1, k_2\) и новое сообщение \(m'\), где \(m'_n = m_n \oplus 0^{127}1\). Мы не знаем промежуточных значений \(c_i\). Чтобы найти \(Tag(m')\), необходимо ``откатиться'' на один шаг назад. Считаем вызовы AES (прямые \(E\) и обратные \(E^{-1}\)):

\begin{enumerate}
\def\labelenumi{\arabic{enumi}.}
\tightlist
\item
  Узнаем последний блок цепи: \(c_n = E^{-1}(k_2, Tag(m))\). \textbf{(1 вызов)}
\item
  Снимаем шифрование \(k_1\), чтобы узнать вход XOR: \(X = E^{-1}(k_1, c_n)\). \textbf{(1 вызов)} \emph{Заметим, что \(X = m_n \oplus c_{n-1}\).}
\item
  Формируем новый вход для последнего блока. Так как \(m'_n = m_n \oplus 0^{127}1\), новый XOR-вход равен: \(X' = m'_n \oplus c_{n-1} = (m_n \oplus 0^{127}1) \oplus c_{n-1} = X \oplus 0^{127}1\). Шифруем его: \(c'_n = E(k_1, X')\). \textbf{(1 вызов)}
\item
  Вычисляем финальный тэг: \(Tag(m') = E(k_2, c'_n)\). \textbf{(1 вызов)}
\end{enumerate}

Итого: \(1 + 1 + 1 + 1 = 4\) вызова функции AES (два дешифрования и два шифрования).

\textbf{Ответ: 4}

\subsection{Блок 5}\label{ux431ux43bux43eux43a-5}

\subsubsection{Задача 1. Доказательство устойчивости композиции хеш-функций}\label{ux437ux430ux434ux430ux447ux430-1.-ux434ux43eux43aux430ux437ux430ux442ux435ux43bux44cux441ux442ux432ux43e-ux443ux441ux442ux43eux439ux447ux438ux432ux43eux441ux442ux438-ux43aux43eux43cux43fux43eux437ux438ux446ux438ux438-ux445ux435ux448-ux444ux443ux43dux43aux446ux438ux439}

\textbf{Решение:} Дано: \(x \neq y\), такие что \(H_2(H_1(x)) = H_2(H_1(y))\). Обозначим промежуточные значения: \(h_x = H_1(x)\) \(h_y = H_1(y)\)

Подставим в исходное равенство: \(H_2(h_x) = H_2(h_y)\).

В этой ситуации логически возможны только два взаимоисключающих случая:

\begin{enumerate}
\def\labelenumi{\arabic{enumi}.}
\tightlist
\item
  \textbf{Случай 1:} \(h_x = h_y\). Возвращаясь к исходным обозначениям, получаем \(H_1(x) = H_1(y)\) при \(x \neq y\). Это означает, что найдена \textbf{коллизия для функции \(H_1\)}.
\item
  \textbf{Случай 2:} \(h_x \neq h_y\). Так как \(H_2(h_x) = H_2(h_y)\) при разных входах, это означает, что найдена \textbf{коллизия для функции \(H_2\)}.
\end{enumerate}

Следовательно, верным должно быть следующее утверждение (или логически эквивалентное ему):

\textbf{Ответ:} Среди предложенных вариантов нужно выбрать тот, который гласит: \textbf{«Либо \(H_1(x) = H_1(y)\) (что дает коллизию для \(H_1\)), либо \(H_1(x) \neq H_1(y)\) (что дает коллизию для \(H_2\))».}

\begin{center}\rule{0.5\linewidth}{0.5pt}\end{center}

\subsubsection{Задача 2. Многосторонняя коллизия (Обобщенный парадокс дней рождений)}\label{ux437ux430ux434ux430ux447ux430-2.-ux43cux43dux43eux433ux43eux441ux442ux43eux440ux43eux43dux43dux44fux44f-ux43aux43eux43bux43bux438ux437ux438ux44f-ux43eux431ux43eux431ux449ux435ux43dux43dux44bux439-ux43fux430ux440ux430ux434ux43eux43aux441-ux434ux43dux435ux439-ux440ux43eux436ux434ux435ux43dux438ux439}

\textbf{Решение:} Обозначим размер выходного пространства \(N = |T|\), а количество случайных выборок \(S\). Нам нужно найти трехстороннюю коллизию (\(k=3\)), то есть \(H(x) = H(y) = H(z)\) для различных \(x, y, z\). Количество возможных троек из \(S\) выборок равно биномиальному коэффициенту: \(\binom{S}{3} \approx \frac{S^3}{6}\). Вероятность того, что конкретная тройка образует коллизию (все три хеша равны конкретному значению из \(N\) возможных) составляет: \(P = \frac{1}{N} \cdot \frac{1}{N} = \frac{1}{N^2}\). Ожидаемое количество трехсторонних коллизий \(E\): \(E \approx \frac{S^3}{6} \cdot \frac{1}{N^2}\). Чтобы найти хотя бы одну коллизию с высокой вероятностью, ожидаемое значение должно быть \(E \ge 1\): \(\frac{S^3}{6N^2} \approx 1 \implies S^3 \approx 6N^2 \implies S \propto (N^2)^{1/3} = N^{2/3}\). Следовательно, асимптотическая сложность поиска трехсторонней коллизии составляет \(O(|T|^{2/3})\).

\textbf{Ответ: Выбираем \(O(|T|^{2/3})\)}

\begin{center}\rule{0.5\linewidth}{0.5pt}\end{center}

\subsubsection{Задача 3. Безопасность настраиваемого (tweakable) блочного шифра}\label{ux437ux430ux434ux430ux447ux430-3.-ux431ux435ux437ux43eux43fux430ux441ux43dux43eux441ux442ux44c-ux43dux430ux441ux442ux440ux430ux438ux432ux430ux435ux43cux43eux433ux43e-tweakable-ux431ux43bux43eux447ux43dux43eux433ux43e-ux448ux438ux444ux440ux430}

\textbf{Решение:} Заданная конструкция: \(E'((k_1,k_2), t, x) = E(k_1, x) \oplus E(k_2, t)\). \textbf{Ответ: Нет, конструкция абсолютно небезопасна.}

\textbf{Доказательство (атака различения):} Настраиваемый блочный шифр должен вести себя как семейство независимых псевдослучайных перестановок для каждого значения \(t\). Построим тривиальный различитель (distinguisher), используя свойство линейности операции XOR.

\begin{enumerate}
\def\labelenumi{\arabic{enumi}.}
\tightlist
\item
  Выберем два произвольных открытых текста \(x_1\) и \(x_2\) (\(x_1 \neq x_2\)).
\item
  Запросим шифротексты для них с использованием некоего tweak \(t_A\): \(C_{1A} = E(k_1, x_1) \oplus E(k_2, t_A)\) \(C_{2A} = E(k_1, x_2) \oplus E(k_2, t_A)\)
\item
  Вычислим их XOR-сумму. Заметим, что член зависящий от tweak сокращается: \(\Delta C_A = C_{1A} \oplus C_{2A} = E(k_1, x_1) \oplus E(k_1, x_2)\)
\item
  Запросим шифротексты для тех же текстов \(x_1\) и \(x_2\), но с \textbf{другим} tweak \(t_B\) (\(t_A \neq t_B\)): \(C_{1B} = E(k_1, x_1) \oplus E(k_2, t_B)\) \(C_{2B} = E(k_1, x_2) \oplus E(k_2, t_B)\)
\item
  Вычислим их XOR-сумму: \(\Delta C_B = C_{1B} \oplus C_{2B} = E(k_1, x_1) \oplus E(k_1, x_2)\)
\end{enumerate}

\textbf{Вывод:} Мы видим, что \(\Delta C_A = \Delta C_B\) выполняется со 100\% вероятностью. В истинном настраиваемом блочном шифре перестановки для \(t_A\) и \(t_B\) независимы, и вероятность совпадения таких разностей пренебрежимо мала (\(\approx 1/2^n\)). Различитель успешно взламывает конструкцию за 4 запроса.

\end{multicols*}
\end{document}

